\section{\ARESBench{} Evaluation Framework}
\label{sec:aresbench}

To validate the reliability model empirically, we built
\ARESBench{}: an open-source, Docker-deployable testbed for measuring
AI security agent reliability under controlled conditions. The design
prioritises reproducibility (single-command deployment, all data
sources public and versioned), generalisability (multiple agent
architectures and threat classes), and controllability (systematic
mismatch injection across a structured matrix). The full architecture
diagram, data-source manifest, and measurement-pipeline details are
maintained in the project artefact repository (see Section~\ref{sec:conclusion}).

\textbf{Agents under test.} \ARESBench{} evaluates four agent
architectures spanning the major paradigms: \textbf{AUT-1
(LangGraph)}~\cite{langgraph2024}, a stateful graph-based ReAct loop;
\textbf{AUT-2 (AutoGen)}~\cite{wu2023autogen}, a three-stage
multi-agent pipeline (triage $\to$ enrichment $\to$ recommendation);
\textbf{AUT-3 (CrewAI)}, a role-based crew (analyst, enrichment,
incident commander); and \textbf{AUT-4}, a deterministic Sigma/YARA
evaluator with no LLM component, serving as the $\delta_R = 0$
control. Hallucination is structurally impossible in AUT-4. In live mode the
testbed drives these agents with a local Mistral-7B-Instruct backbone
via Ollama; the released \SAFK{} corpus, however, uses the analytical
path (below), and the live calibration and cascade studies use hosted
frontier models, not Ollama.

\textbf{Controlled mismatch injection.} A Mismatch Injector applies
controlled degradation across four dimensions ($\delta_S, \delta_K,
\delta_C, \delta_R$) at four levels each (0, 0.33, 0.66, 1.0).
Mechanisms are deliberately realistic: disabling tool nodes in the
agent's execution graph ($\delta_S$); filtering ATT\&CK TTP context
from retrieval to retain only a controlled fraction ($\delta_K$);
intercepting API calls via a middleware proxy that returns
permission-denied errors ($\delta_C$); and inserting misleading content
into the system prompt at controlled probability ($\delta_R$).
Temporal drift ($\delta_T$) is set parametrically via the agent's
knowledge-age attribute.

\textbf{Failure measurement.} An Anomaly Monitor parses agent
output, cross-references claims against ATT\&CK and NVD as a
closed-world ground truth, and flags any unverifiable claim asserted
with confidence $> 0.70$ as a hallucination candidate. It then
computes the five mismatch dimensions, the ARS, and the Epistemic Gap,
and assigns taxonomy labels using a rule-based engine validated
against human annotations (Cohen's $\kappa = 0.84$ for origin,
$\kappa = 0.91$ for manifestation).

\textbf{From injection to measurement.} Controlled injection is how we
\textit{construct} the study, not how the model is instantiated in
deployment: each dimension is computed from artefacts a SOC already
possesses. \textit{Skill} ($\delta_S$) and \textit{capability}
($\delta_C$) are read from the agent's registered tool and function
manifest, the fraction of the task's required procedures and APIs it is
not wired to invoke. \textit{Knowledge} ($\delta_K$) is measured over
context, not weights: the fraction of the task kill-chain's ATT\&CK TTP
nodes absent from the agent's retrieval index or supplied context,
exactly computable for a retrieval-augmented agent; a genuinely novel
threat matches nothing and drives $\delta_K \to 1$, precisely the
high-risk regime the score should flag. \textit{Temporal drift}
($\delta_T$) uses time since the last knowledge refresh, and expressed
uncertainty $U$ is the agent's stated confidence. Four of the five
dimensions are thus prospective, needing no run-time ground truth; only
$\delta_R$ requires post-hoc comparison. The live-incident harness
(Empirical Insights) exercises this path on real CVEs, with no
injection.

\textbf{The \SAFK{} corpus (analytical).} \SAFK{} is a labelled corpus
of 1{,}000 traces stratified across five threat classes (200 each),
four agent architectures (250 each), and 16 mismatch conditions.
Crucially, its \textit{outcomes}, hallucination, expressed confidence,
task success, and the taxonomy labels, are generated from the mismatch
inputs by calibrated closed-form models rather than read from live
model execution. This makes \SAFK{} an analytical instrument: mismatch--failure
relationships can be studied under exact ground truth and reproduced
deterministically, at the cost of encoding rather than discovering
them, which is why this article's central claim is validated not on
\SAFK{} but on live agents. The corpus is released under CC BY 4.0
alongside the code.

\textbf{The live-agent harness.} To test the model's assumptions
against reality, \ARESBench{} also drives \textit{real} agents: it
applies the same $\delta_K$ manipulation to the prompt of a
frontier LLM, then measures hallucination against ATT\&CK ground truth
and records the agent's own expressed confidence. This is the path
behind the calibration-gap result and the live-incident study
(Empirical Insights); the harness and its 600-execution output are
released with the corpus so the comparison is reproducible. Because
frontier models are stochastic and served behind evolving endpoints,
re-running the harness reproduces the reported \textit{effects}
(rising uncertainty and flat fabrication under increasing $\delta_K$)
directionally rather than the exact per-call values; the committed
run log is the frozen record behind every number we report.
